\chapter{Unity Simulation Design}\label{ch:unity}
\thispagestyle{pagestyle}

The simulation is not just a backdrop -- it is the only reality the
agent ever experiences during training. Every aspect of the environment,
from how the car steers around a corner to what penalty it receives for
scraping a wall, directly shapes the policy that eventually runs on the
real vehicle. This chapter explains the design choices made for the Unity
simulation and why they were made that way.

\section{Scene Construction}
\label{sec:scene}

The simulation scene is constructed programmatically at runtime rather
than authored in the Unity Editor. \texttt{RuntimeSceneBuilder.cs} runs
on scene load and checks whether the required game objects are present;
if absent, it invokes \texttt{CarFactory} to spawn the full vehicle
hierarchy and \texttt{SpawnObstacleCourse} to instantiate the track. This
approach guarantees a well-defined initial state regardless of how the
Editor project was last saved.

The obstacle course is a circular ring generated by
\texttt{ObstacleCourse.cs}. The ring is bounded by two concentric
cylindrical walls: inner radius 25~m, outer radius 40~m, giving a track
width of 15~m. Twelve obstacle objects are distributed at equal angular
intervals. Wall height is 3~m, sufficient to be detected by raycasts but
low enough to avoid occluding the follow camera. The vehicle is spawned
at position $(32.5,\ 0.5,\ 0)$ --- on the track centreline facing the
initial direction of travel.

\section{Vehicle Physics Model}
\label{sec:vehicle-physics}

The vehicle is a rear-wheel-drive configuration assembled by
\texttt{CarFactory.cs}. The chassis mass is 1500~kg. Four
\texttt{WheelCollider} components provide the physics contact model using
Unity's \cite{UnityManual2022} built-in Pacejka-inspired tyre model.
Key parameters are listed in \cref{table:vehicle-params}.

\begin{table}[H]
\caption{Key vehicle physics parameters.}
\label{table:vehicle-params}
\begin{tabular}{|p{5cm}|p{2.5cm}|p{2cm}|}
  \hline
  \textbf{Parameter} & \textbf{Value} & \textbf{Unit} \\
  \hline
  Chassis mass           & 1500   & kg \\
  \hline
  Maximum motor torque   & 450    & Nm \\
  \hline
  Maximum steer angle    & 35     & degrees \\
  \hline
  Final drive ratio      & 9.0    & -- \\
  \hline
  Maximum brake torque   & 3000   & Nm \\
  \hline
  Regenerative brake torque & 150 & Nm \\
  \hline
  Maximum speed          & 150    & km/h \\
  \hline
\end{tabular}
\end{table}

\texttt{VehicleController.cs} translates the three-component action
vector $[a_{\mathrm{steer}}, a_{\mathrm{throttle}}, a_{\mathrm{brake}}]$
received from \texttt{RemoteDriveInput.cs} into WheelCollider motor
torques and steer angles. A creep behaviour applies a small constant
torque in drive gear when both throttle and brake inputs are near zero,
mirroring a real automatic-transmission vehicle.

\section{Perception Layer}
\label{sec:perception}

The vehicle's sole sensor is a LiDAR-inspired raycast fan implemented in
\texttt{RaycastSensor.cs}. Twenty-one rays are cast in the vehicle's
horizontal plane, spanning a $120^\circ$ arc centred on the forward axis,
with a range of 50~m. The raycast loop itself is compact:

\begin{lstlisting}[language={[Sharp]C},
  caption={Ray fan loop from RaycastSensor.cs.},
  label={code:raycast}]
void CastRays()
{
    Vector3 origin = transform.TransformPoint(rayOriginOffset);
    float halfSpread = raySpreadAngle * 0.5f;
    float step = rayCount > 1 ? raySpreadAngle / (rayCount - 1) : 0f;

    for (int i = 0; i < rayCount; i++)
    {
        float angle = -halfSpread + step * i;
        Vector3 direction = Quaternion.Euler(0f, angle, 0f) * transform.forward;

        if (Physics.Raycast(origin, direction, out RaycastHit hit, rayLength, hitLayers))
            _distances[i] = hit.distance / rayLength; // 0 = close, 1 = max range
        else
            _distances[i] = 1f;                       // no obstacle detected
    }
}
\end{lstlisting}

Dividing by \texttt{rayLength} normalises each reading into $[0, 1]$:
a value of 1 means the ray reached full range without hitting anything
while values near 0 indicate an obstacle very close to the bumper.
Each ray returns a normalised distance $d_i \in [0, 1]$, and the sensor
publishes the full array to \texttt{env0/vehicle/training/obs} at 10~Hz.

The complete observation vector received by Python is:
\begin{equation}
  \mathbf{o} = \left[ d_0,\, \ldots,\, d_{N-1},\,
                      \frac{v}{v_{\max}},\,
                      \frac{\delta + 1}{2} \right] \in \mathbb{R}^{N+2},
  \label{eq:obs}
\end{equation}
where $N$ is the configured ray count, $v$ is current speed (km/h),
$v_{\max} = 150$~km/h, and $\delta$ is steer angle normalised to
$[-1, 1]$. All elements of $\mathbf{o}$ lie in $[0, 1]$, simplifying
network initialisation.

\section{Reward Function}
\label{sec:reward}

The per-step reward is:
\begin{equation}
  r_t = r_{\mathrm{progress}} + r_{\mathrm{collision}} + r_{\mathrm{milestone}},
  \label{eq:reward}
\end{equation}
where:
\begin{itemize}[topsep=1.15pt,itemsep=1.15pt]
  \item $r_{\mathrm{progress}} = 2 \cdot \Delta p_t$ --- signed change in
    track angle progress (degrees) since the last step.
  \item $r_{\mathrm{collision}} = -5$ on any step where the vehicle
    reports a collision contact (applied once per event).
  \item $r_{\mathrm{milestone}} = +50$ on full lap completion;
    $+25$ on first half-turn of the course (at most once per episode
    each).
\end{itemize}

The reward shaping is implemented in the \texttt{FixedUpdate()} loop of
\texttt{TrainingBridge.cs}. Collision detection uses \texttt{OnCollisionEnter}
rather than the raycast array, preventing any edge-case where a collision
could be missed if the sensor update and physics update happened to
occur in the wrong order:

\begin{lstlisting}[language={[Sharp]C},
  caption={Reward computation excerpt from TrainingBridge.cs.},
  label={code:reward}]
float angleDelta = Mathf.DeltaAngle(_lastAngle, currentAngle);
_lastAngle = currentAngle;
_totalAngleProgress += angleDelta;

float stepReward = angleDelta * progressReward;   // progressReward = 2

bool collided = _collidedThisStep;
_collidedThisStep = false;
if (collided)
    stepReward += collisionPenalty;               // collisionPenalty = -5

if (!_halfTurnAwarded)
{
    float heading = vehicle.transform.eulerAngles.y;
    if (heading >= 170f && heading <= 190f)
    {
        stepReward += halfTurnBonus;              // halfTurnBonus = 25
        _halfTurnAwarded = true;
    }
}

if (obstacleCourse != null && _totalAngleProgress >= 360f)
{
    stepReward += finishBonus;                    // finishBonus = 50
    finished = true;
}
\end{lstlisting}

An episode terminates on lap completion, collision, or after the
60-second wall-clock timeout. Timeout terminations are reported as
\texttt{truncated=True} so that Stable-Baselines3 bootstraps the value
estimate rather than treating the transition as a terminal failure.

\section{MQTT Bridge in Unity}
\label{sec:unity-mqtt}

Unity does not ship a production-grade MQTT client. SHILATE includes a
minimal custom \texttt{MqttClient.cs} that implements the MQTT 3.1.1 wire
protocol \cite{OASIS2014} natively in C\# without external dependencies.
The client handles \texttt{CONNECT}, \texttt{PUBLISH},
\texttt{SUBSCRIBE}, \texttt{PINGREQ/PINGRESP}, and \texttt{DISCONNECT}
packet types with a 30-second keep-alive.

\texttt{LedaBroker.cs} provides the application-facing API, configured
via \texttt{LedaBroker.Configure(host, port, "env0")} at startup. It
publishes vehicle telemetry at 10~Hz and receives control commands. A
Python heartbeat timeout (5-second threshold) detects when the trainer
has stopped sending commands and holds the vehicle stationary until
reconnection. This is a small but important detail --- without it, a
crashed Python process would leave the simulated car coasting forward
indefinitely into a wall.

\section{Training Controller Editor Window}
\label{sec:editor-window}

The \textit{SHILATE Training Controller}
(\texttt{SHILATE} $\to$ \texttt{Training Controller}) is an EditorWindow
that eliminates manual terminal management during training. Its layout
from top to bottom comprises:
\begin{enumerate}[leftmargin=2cm,topsep=1.15pt,itemsep=1.15pt,partopsep=1.15pt,parsep=1.15pt]
  \item \textbf{Toolbar}: training duration and MQTT connection indicator.
  \item \textbf{Health banner}: colour-coded strip showing the current
    state (\textit{Idle}, \textit{Healthy}, \textit{Warning},
    \textit{Critical}, \textit{Disconnected}) with explicit problem text.
  \item \textbf{Snapshot row}: episode count, last-episode reward,
    10-episode rolling mean, and heartbeat counter.
  \item \textbf{Settings panel} (collapsible): all TrainingSettings
    fields exposed as GUI controls.
  \item \textbf{Metric graphs}: scrolling plots of episode reward mean,
    policy gradient loss, value loss, and KL divergence.
  \item \textbf{Issue log}: timestamped list of every health-state
    transition for the session.
  \item \textbf{Controls}: \textit{Start Training} / \textit{Run Model}
    / \textit{Stop}.
\end{enumerate}
The \textit{Start Training} button enters Play mode programmatically and,
once the scene is loaded and the MQTT connection is confirmed, launches
the Python training script with arguments from TrainingSettings. The
health evaluator monitors five independent failure modes: MQTT
disconnection, Python process crash, observation gap exceeding 5~seconds,
heartbeat/reward gap exceeding 60~seconds, and reward collapse (more than
50\% drop over a 10-episode window).
